% !TEX program = xelatex
\documentclass[11pt,a4paper]{article}

\usepackage[margin=1.8cm]{geometry}
\usepackage{fontspec}
\usepackage{polyglossia}
\setmainlanguage{spanish}
\setmainfont{Latin Modern Roman}
\usepackage{amsmath,amssymb}
\usepackage{enumitem}
\usepackage{fancyhdr}
\usepackage[hidelinks]{hyperref}
\usepackage{parskip}

\setlength{\headheight}{14pt}
\setlength{\headsep}{0.45cm}
\pagestyle{fancy}
\fancyhf{}
\fancyhead[L]{\small Grupo de Estudio: Teoría Analítica de Números}
\fancyhead[R]{\small\today}
\renewcommand{\headrulewidth}{0.4pt}
\fancypagestyle{plain}{%
  \fancyhf{}
  \fancyhead[L]{\small Grupo de Estudio: Teoría Analítica de Números}
  \fancyhead[R]{\small\today}
  \renewcommand{\headrulewidth}{0.4pt}
}

\newcommand{\PP}{\mathcal{P}}
\newcommand{\card}[1]{\lvert #1\rvert}
\newcommand{\Set}[1]{\left\{#1\right\}}

\title{Tarea de la Clase 1\\\large Teoría de conjuntos}
\author{}
\date{}

\begin{document}
\maketitle
\vspace{-1.2cm}
\noindent{\small Autor: Daniel Soto Villada}\par
\vspace{0.35cm}
\noindent Nombre: \rule{8cm}{0.4pt}\hfill Fecha: \rule{3cm}{0.4pt}

\noindent\textbf{1.} Relacione cada concepto de la columna izquierda con un único significado de la columna derecha: trace una línea entre ellos. Hay \textbf{dos opciones sobrantes}; cada significado correcto se usa una sola vez.

\vspace{0.2cm}
\small
\noindent
\begin{minipage}[t]{0.28\textwidth}
\textbf{Conceptos}
\begin{enumerate}[label=\Alph*., leftmargin=*, itemsep=0.95em, topsep=0.4em]
    \item $x\in A$
    \item $A\subseteq B$
    \item $A=B$
    \item $\varnothing\subseteq A$
    \item $\PP(A)$
    \item $A\cup B$
    \item $A\cap B$
    \item $A\setminus B$
    \item $A^c=X\setminus A$
    \item $(a,b)$
    \item $A\times B$
    \item $\card{A\times B}=\card{A}\,\card{B}$
\end{enumerate}
\end{minipage}\hfill
\begin{minipage}[t]{0.61\textwidth}
\textbf{Significados}
\begin{enumerate}[label=\arabic*., leftmargin=*, itemsep=0.38em, topsep=0.4em]
    \item Reúne los elementos que están en ambos conjuntos; pertenecer a uno solo no basta.
    \item Si los conjuntos son finitos, el número de pares que se forman es el producto del número de opciones para cada coordenada.
    \item Conjunto de todos los subconjuntos de $A$, incluidos $\varnothing$ y $A$.
    \item Afirma que $x$ es uno de los elementos del conjunto $A$.
    \item Conjunto de los elementos de $A$ que no pertenecen a $B$; en general, invertir los conjuntos cambia el resultado.
    \item Dos pares son iguales exactamente cuando coinciden sus primeras coordenadas y también sus segundas; cambiar el orden puede cambiar el par.
    \item Conjunto de pares cuya primera coordenada se elige en $A$ y la segunda en $B$; intercambiar los factores puede cambiarlo.
    \item Conjunto de los elementos de $X$ que no pertenecen a $A$; su definición depende del universo fijado.
    \item Se cumple cuando todo elemento del primero también pertenece al segundo; compara conjuntos, no un elemento con un conjunto.
    \item Contiene los elementos que pertenecen al primero \emph{o} al segundo, incluyendo los que pertenecen a ambos.
    \item Se cumple si cada conjunto está incluido en el otro; tener igual número de elementos no es, por sí solo, suficiente.
    \item Es subconjunto de cualquier conjunto, pues no existe un elemento suyo que pueda incumplir la condición de inclusión.
    \item Basta con que los dos conjuntos tengan la misma cantidad de elementos, aunque sus elementos sean distintos.
    \item Es elemento de todo conjunto, porque es subconjunto de todo conjunto.
\end{enumerate}
\end{minipage}

\noindent\textbf{Reto.} Identifique las dos opciones sobrantes y explique brevemente qué confusión representa cada una.

\noindent Para los puntos 2 y 3, seleccione una sola respuesta. Las opciones se presentan por comprensión.

\noindent\textbf{2.} El universo está formado por los enteros desde cero hasta doce, incluyendo ambos extremos. El primer conjunto contiene los números pares de ese universo, y el segundo contiene los enteros estrictamente mayores que cuatro. ¿Cuál opción describe los elementos del primer conjunto que no pertenecen al segundo?

\begin{enumerate}[label=\Alph*), leftmargin=1.2cm, itemsep=0.45em]
    \item $\Set{x\in\mathbb{Z}:0\leq x\leq12,\ x\text{ es par},\ x\leq4}$
    \item $\Set{x\in\mathbb{Z}:0\leq x\leq12,\ x\text{ es par},\ x<4}$
    \item $\Set{x\in\mathbb{Z}:0\leq x\leq12,\ x\text{ es par o }x\leq4}$
    \item $\Set{x\in\mathbb{Z}:0\leq x\leq12,\ x\text{ es par},\ x>4}$
\end{enumerate}

\noindent\textbf{3.} El universo está formado por todos los enteros desde menos ocho hasta ocho, incluyendo ambos extremos. El primer conjunto contiene los múltiplos de tres de ese universo; el segundo contiene los enteros pares positivos. ¿Cuál opción describe exactamente los elementos del universo que no pertenecen a ninguno de los dos conjuntos?

\begin{enumerate}[label=\Alph*), leftmargin=1.2cm, itemsep=0.45em]
    \item $\Set{x\in\mathbb{Z}:-8\leq x\leq8,\ x\text{ no es múltiplo de }3,\ (x\text{ es impar o }x\leq0)}$
    \item $\Set{x\in\mathbb{Z}:-8\leq x\leq8,\ (x\text{ no es múltiplo de }3)\text{ o }(x\text{ es impar o }x\leq0)}$
    \item $\Set{x\in\mathbb{Z}:-8\leq x\leq8,\ (x\text{ es múltiplo de }3)\text{ o }(x\text{ es par positivo})}$
    \item $\Set{x\in\mathbb{Z}:-8\leq x\leq8,\ x\text{ no es múltiplo de }3,\ x\text{ es impar},\ x\leq0}$
\end{enumerate}


\noindent En los puntos 4 y 5, localice el primer razonamiento inválido. No basta con señalar que la conclusión es falsa: explique el error y responda lo que se pide.

\noindent\textbf{4.} Un estudiante afirma que la siguiente identidad vale para cualesquiera conjuntos $A$, $B$ y $C$:
\[
A\setminus(B\cup C)=(A\setminus B)\cup(A\setminus C).
\]
Para justificarla, escribe, para un elemento arbitrario $x$:
\begin{align*}
x\in A\setminus(B\cup C)
&\Longleftrightarrow x\in A\text{ y }x\notin B\cup C\\
&\Longleftrightarrow x\in A\text{ y }(x\notin B\text{ o }x\notin C)\\
&\Longleftrightarrow (x\in A\text{ y }x\notin B)\text{ o }(x\in A\text{ y }x\notin C).
\end{align*}
Indique cuál equivalencia falla y cuál conectiva debe aparecer en su lugar. Después, dé conjuntos finitos $A$, $B$ y $C$ para los que los dos lados de la identidad propuesta sean distintos, calcule ambos lados y escriba la identidad correcta.

\noindent\textbf{5.} Otro estudiante sostiene que, para cualesquiera conjuntos $A$ y $B$,
\[
\PP(A\cup B)=\PP(A)\cup\PP(B).
\]
Su argumento comienza así: «Sea $S\in\PP(A\cup B)$. Entonces cada elemento de $S$ pertenece a $A$ o a $B$. Por consiguiente, $S\subseteq A$ o $S\subseteq B$».
Señale el salto lógico injustificado. Construya conjuntos finitos $A$ y $B$ y un conjunto $S$ que refuten la igualdad. Finalmente, demuestre cuál de las dos inclusiones entre $\PP(A\cup B)$ y $\PP(A)\cup\PP(B)$ sí se cumple siempre.

\noindent En los puntos 6 y 7, escriba demostraciones completas y justifique las inclusiones o igualdades elemento por elemento.

\noindent\textbf{6.} Demuestre que, para cualesquiera conjuntos $A$ y $B$,
\[
A\subseteq B\quad\Longleftrightarrow\quad\PP(A)\subseteq\PP(B).
\]
Pruebe ambas direcciones. En la recíproca, explique cómo pasar de la inclusión entre conjuntos potencia a una afirmación sobre un elemento arbitrario de $A$.
\emph{Pista.} Para cada $x\in A$, considere el conjunto unitario $\{x\}$ y pregúntese en qué conjunto potencia pertenece.

\noindent\textbf{7.} Demuestre la siguiente caracterización, que muestra cuándo se puede intercambiar el orden de los factores sin cambiar el producto cartesiano:
\[
A\times B=B\times A
\quad\Longleftrightarrow\quad
A=B\text{ o }A=\varnothing\text{ o }B=\varnothing.
\]
Trate por separado los casos en que alguno de los conjuntos es vacío y aquellos en que ambos son no vacíos. En estos últimos, use la igualdad de pares ordenados para justificar las dos inclusiones entre $A$ y $B$.
\emph{Pista.} Si ambos conjuntos son no vacíos, fije un elemento de uno de ellos y forme pares ordenados cuya primera coordenada recorra el otro conjunto. Use la igualdad de productos para determinar dónde debe pertenecer esa primera coordenada.

\end{document}
